% =============================================================================
\section{Introduction}
\label{sec:intro}
% =============================================================================

Over the past decade, membership inference attacks
(MIA)~\cite{homer2008resolving,shokri2017membership} have become the
standard tool for measuring how much a trained model leaks about its
training data.
Hundreds of papers have proposed stronger
attacks~\cite{salem2019ml,nasr2019comprehensive,choquette2021label,song2021systematic,carlini2022lira,ye2022enhanced,watson2022importance,zarifzadeh2024rmia},
dedicated defenses~\cite{nasr2018machine,jia2019memguard,abadi2016deep},
tighter theoretical
bounds~\cite{yeom2018privacy,sablayrolles2019whitebox}, and
systematizations of the resulting
landscape~\cite{papernot2018sok,hu2022membership}.
The same model-centric view organizes the sibling lines of work:
attribute and property
inference~\cite{fredrikson2015model,ganju2018property}, training-data
extraction~\cite{carlini2019secret,carlini2021extracting}, differential
privacy~\cite{dwork2006calibrating,abadi2016deep}, and machine
unlearning~\cite{cao2015towards,bourtoule2021machine,chen2021unlearning}
all quantify or remove leakage from a trained model.
Yet a more fundamental question has gone largely unanswered:

\smallskip
\begin{center}
Why do some datasets leak more than others,\\
independent of the model trained on them?
\end{center}
\smallskip

This is not only an engineering question.
It is also a question about the nature of data itself.
The pre-ML privacy literature established that re-identification risk is
driven by intrinsic properties of the data: the sparsity of the Netflix
ratings matrix~\cite{narayanan2008robust}, the uniqueness of mobility
traces~\cite{demontjoye2013unique} and credit-card
metadata~\cite{demontjoye2015unique}, and the incompleteness that
defeats naive anonymization~\cite{sweeney2002k,rocher2019estimating}.
Machine learning privacy has largely set this lesson aside, treating the
dataset as a fixed backdrop behind the model. We argue that membership
inference is, in part, a symptom of how the training data is shaped.

\paragraph{The problem with current practice.}
Every MIA paper chooses a dataset, trains a model, and reports an AUC.
The AUC is then attributed to the attack method or the defense mechanism.
But when we decompose the variation in AUC across 31 datasets, three
attacks, and three model families, the dataset is anything but a fixed
backdrop. Among the regularized models that practitioners actually deploy,
dataset choice explains $49.5\%$ of the variation in AUC, twice the share
of model architecture ($24.4\%$) and well above the attack method
($26.1\%$). Only a deliberately overfit model, an unbounded-depth Random
Forest that memorizes its training set outright, erases this margin
(Section~\ref{sec:finding1}).
Dataset structure is a first-class driver of privacy risk,
yet it is rarely reported or controlled for.

A second pathology is benchmark reliance.
Purchase100 and Texas100~\cite{shokri2017membership} are the two most
common MIA benchmarks.
We show that both are geometric outliers. Their sample uniqueness is the
highest among all 31 datasets we study, placing them far outside the range
of typical deployed data, so papers that report results only on these
benchmarks measure privacy risk in an unrepresentatively high-risk regime.

\paragraph{Our approach.}
We ask a temporally prior question:
Can the privacy risk of a dataset be anticipated before any model is trained?
If so, practitioners gain an actionable signal: estimate risk at data
collection time, before committing to a training pipeline or a privacy budget.

We give a qualified yes. We introduce the Dataset Privacy Risk Index (DPRI),
a scalar pre-training risk estimator derived from distributional statistics:
sample uniqueness, local density, outlier score, feature entropy, class
cluster separation, and data dimensionality.
Each statistic is motivated by a theoretical result
(Theorem~\ref{thm:main}) showing that MIA advantage is bounded by a
function of local data geometry.
Across 31 datasets, DPRI ranks datasets by risk with a Spearman correlation
of $0.61$ ($p < 0.001$). We report this number under nested cross-validation
that selects features on training folds only, so it is free of the selection
optimism that inflates small-$n$ results; the same discipline is what
exposed our own early estimates, on fewer datasets, as artifacts.
Data geometry significantly shapes privacy risk without fully determining
it.

Our question is related to, but distinct from, the memorization literature.
Feldman~\cite{feldman2020memorization} showed that rare individual samples
must be memorized for good generalization, a phenomenon observable only
after training. We ask the complementary and temporally prior question:
can dataset-level structural statistics predict aggregate membership
inference risk before any model is trained? Where that line of work
identifies which samples a given model will memorize (post-training,
sample-level), we predict how much a dataset will leak (pre-training,
dataset-level), a distinction that is both scientific and operational. A
pre-training score is the only kind a practitioner can act on at data
collection time.

\paragraph{Contributions.}
\begin{itemize}

\item \textbf{Theoretical foundation.}
We prove that the Bayes-optimal MIA advantage for a sample $x$ is
upper-bounded by $O(u(x) / \rho(x)^{1/d})$, where $u(x)$ is sample
uniqueness and $\rho(x)$ is local density, both computable from raw data
(Theorem~\ref{thm:main}).
This provides the first formal link between pre-training data geometry
and membership inference risk.

\item \textbf{Dataset Privacy Risk Index (DPRI).}
We define a dataset-level risk score from six distributional features and
validate it across 31 datasets in five domains using three MIA attacks and
three model families. DPRI ranks datasets by risk with an unbiased nested
cross-validated Spearman correlation of $0.61$ ($p < 0.001$).

\item \textbf{A calibrated, qualified finding.}
The prediction is moderate, not precise, and no single statistic suffices.
Sample uniqueness, local density, class separability, and dimensionality
are each significant predictors, yet privacy risk is shaped by data geometry
without being determined by it. We argue this nuance is itself a
contribution. It bounds what any pre-training predictor can achieve.

\item \textbf{Two measurement findings.}
(F1)~Among realistic, regularized models, dataset structure governs MIA
risk more than model architecture ($49.5\%$ vs.\ $24.4\%$ of AUC variation)
and far more than the choice of attack ($26.1\%$), so an attack number
reported on a single benchmark conflates three distinct sources of
variation, and the dataset is the largest of them.
(F2)~The benchmarks that anchor a decade of MIA research are geometric
outliers, casting doubt on the external validity of results reported only
on them.

\item \textbf{Open-source tool.}
We release DPRI as a Python package that audits any tabular dataset in
minutes on a laptop, with no model training and no GPU, outputting a risk
score and its distributional breakdown.

\end{itemize}

\paragraph{Scope.}
We study 31 public datasets spanning tabular, image, network, medical, and
recommendation domains, in the centralized training setting.
Extensions to federated learning and generative models are left as future
work.
All datasets used are publicly available. No proprietary data was accessed.
